\documentclass[11pt,a4paper]{article}
\usepackage[margin=1in]{geometry}
\usepackage{booktabs}
\usepackage{xcolor}
\usepackage{hyperref}
% enumitem not available on this system
\usepackage{amsmath}
\usepackage{caption}

\definecolor{crit}{HTML}{cf222e}
\definecolor{mod}{HTML}{bf8700}
\definecolor{lightgray}{HTML}{f6f8fa}

\hypersetup{colorlinks=true,linkcolor=black,urlcolor=blue}

\setlength{\parindent}{0pt}
\setlength{\parskip}{6pt}

\title{\textbf{Criticality v2 --- Sequence Scoring Diagnosis}}
\author{SLM Agentic Benchmarking}
\date{February 7, 2026}

\begin{document}
\maketitle
\thispagestyle{empty}

\vspace{-1em}
\begin{center}
\small 200 tasks per model\enspace$\cdot$\enspace 4-choice MCQ\enspace$\cdot$\enspace random baseline = 25\%\\[2pt]
\small Scoring: per-option continuation logprob --- $P(\text{arg}_i \mid \text{``Topic: \{topic\}. The strongest argument is: ''})$
\end{center}

%----------------------------------------------------------------------
\section*{Verdict}

\colorbox{crit!10}{\parbox{\dimexpr\textwidth-2\fboxsep}{%
\textcolor{crit}{\textbf{Benchmark broken.}}
All 9 models cluster within 4.5pp (29.0--33.5\%). No model pair is statistically distinguishable.
The benchmark has zero discriminative power.}}

%----------------------------------------------------------------------
\section{Results}

\begin{table}[h]
\centering
\caption{Per-option continuation scoring. All McNemar pairwise $p > 0.24$.}
\vspace{4pt}
\begin{tabular}{@{}lrccc@{}}
\toprule
\textbf{Model} & \textbf{Acc.} & \textbf{$p$ vs.\ 25\%} & \textbf{Rank Corr.} & \textbf{Cal.\ Error} \\
\midrule
Gemma3-1B           & 33.5\% & 0.004 * & 0.269 & 0.026 \\
Falcon-H1-90M       & 32.0\% & 0.015 * & 0.171 & 0.025 \\
GPT-OSS-20B         & 31.5\% & 0.023 * & 0.273 & 0.014 \\
Gemma3n-E4B         & 31.5\% & 0.023 * & 0.269 & 0.011 \\
Qwen3-0.6B          & 31.0\% & 0.032 * & 0.184 & 0.020 \\
Gemma3-4B           & 30.0\% & 0.063   & 0.276 & 0.003 \\
DASD-4B             & 29.5\% & 0.084   & 0.253 & 0.008 \\
Gemma3n-E2B         & 29.5\% & 0.084   & 0.233 & 0.006 \\
Phi4-Mini-Reasoning & 29.0\% & 0.112   & 0.239 & 0.009 \\
\bottomrule
\end{tabular}
\label{tab:old-results}
\end{table}

%----------------------------------------------------------------------
\section{Problems identified}

\subsection{\textcolor{crit}{Critical:} Ambiguous tasks}

46\% of tasks have a top-2 quality gap below 0.10 on a 0--1 scale. The ``correct'' answer is functionally a coin flip in nearly half the tasks.

\begin{table}[h]
\centering
\small
\begin{tabular}{@{}lc@{}}
\toprule
Gap threshold & \% of tasks below \\
\midrule
$< 0.01$ & 12\% \\
$< 0.05$ & 26\% \\
$< 0.10$ & 46\% \\
$< 0.20$ & 80\% \\
\bottomrule
\end{tabular}
\end{table}

\textbf{Cause:} The dataset contains 758 strong / 203 medium / 39 weak arguments. The task generator fills all 4 MCQ slots from the largest available tier, producing tasks where every option is ``strong'' with near-identical quality scores.

\subsection{\textcolor{crit}{Critical:} Scoring measures fluency, not quality}

Continuation scoring computes how ``expected'' each argument is as natural language. Each option gets a separate \texttt{model.eval()} call---the model never sees the alternatives.

\begin{table}[h]
\centering
\small
\begin{tabular}{@{}lc@{}}
\toprule
Correlation & $r$ \\
\midrule
Quality vs.\ logprob & 0.272 \\
Length vs.\ logprob  & $-0.048$ \\
Length vs.\ quality  & 0.317 \\
\bottomrule
\end{tabular}
\end{table}

The quality--logprob correlation (0.272) is weak and likely driven by the prefix mentioning ``strongest argument'' rather than genuine quality assessment.
High-quality arguments use more specific vocabulary, which \emph{lowers} per-token logprob.

\subsection{\textcolor{crit}{Critical:} No pairwise discrimination}

All McNemar pairwise $p$-values exceed 0.24. The accuracy spread is only 4.5pp, and all 95\% CIs overlap completely. Inter-model agreement ranges from 54--80\%, suggesting models are producing near-identical predictions.

\subsection{\textcolor{mod}{Moderate:} Uniform position bias}

All models prefer option A (28--36\% of predictions vs.\ 26.5\% true base rate). This shared, uniform bias indicates models exploit a positional artifact in the continuation scoring rather than reading argument content.

%----------------------------------------------------------------------
\section{Root cause}

The two critical issues compound:

\begin{enumerate}
\item \textbf{Data:} Nearly half the tasks are ambiguous (gap $<$ 0.10), so even a perfect evaluator would score $\sim$70\% on the ``easy'' subset and $\sim$50\% on the ``hard'' subset---blending to $\sim$35\% overall.
\item \textbf{Method:} Continuation scoring is orthogonal to quality assessment. It measures text probability, which is dominated by syntax and vocabulary frequency, not argumentative strength.
\end{enumerate}

Together, the benchmark asks an ill-posed question (ambiguous tasks) using the wrong tool (fluency scoring). The result is indistinguishable noise.

%----------------------------------------------------------------------
\section{Fixes applied}

\begin{enumerate}
\item \textbf{Task generation (data fix):} Enforce min quality gap $\geq$ 0.15 between the best and second-best argument. Load 4{,}000+ arguments to ensure cross-tier diversity. Result: mean gap = 0.212, 100\% of tasks above threshold.

\item \textbf{Single-token MCQ scoring (method fix):} Present all 4 labeled arguments in one prompt, extract logprobs for A/B/C/D at the answer position. The model now performs comparative judgment in a single forward pass.
\end{enumerate}

These fixes produced an accuracy spread of 29pp (25--54\%), with 17/36 pairwise comparisons statistically significant and 8/9 models clearly above random at $p < 0.001$.

\end{document}
